\section{Introduction}\label{sec:intro}

Input-output models are an established tool for estimating the economic effects of structural policy choices related to public investment, industrial policy or labor market trends \citep{BlairMiller}. They are therefore also widely used to estimate the economic effects of additional investment undertaken in the course of a socio-ecological transformation aiming to render social provisioning processes carbon-neutral \citep{Leontief.1970, Hornykewycz.2025, Markandya.2016, Markaki.2013}. However, it is well known that the results obtained from such models depend crucially on the underlying axiomatic setup and, specifically, on closure choices, i.e.\ assumptions about how to close the model \citep{Robinson.2006, Rose.1995}.

Closure choices encode assumptions about the economic situation a policy meets, and so resonate with economics as ``the art of choosing models which are relevant to the contemporary world'' \citep{Keynes.1938}, emphasizing the importance of socio-historical and institutional context \citep{Polanyi.2001, Hodgson.2010}. Moreover, closure choices also reflect traditional divides in economic theory so that, sometimes, old ghosts haunt modern applications, including those in applied socio-ecological transformation research.

A first example of such a divide is the question whether, and how far, output or prices react to a change in final demand (scarcity vs. output-adjustment): it sits at the core of the division between computable general equilibrium models (CGEs) and the linear extrapolations of the quantities produced by the Leontief inverse \citep{Leontief.1936, BlairMiller}. Moreover, a series of other foundational questions come into view when reflecting on paradigmatic cleavages in input--output modeling: choices on production technology (Cobb Douglas vs. Leontief; \citealp{Labini.1995, Klump.2012, Leontief.1936}), financing (external vs. internal; \citealp{BlinderSolow.1973, Hein.2023}), policy implementation (incentivizing private actors vs. state-led action) or technological change (exogenous vs. endogenous; \citealp{Kaldor1961}) are apt examples, where disagreement can drive a structural wedge into received estimates and projections. To substantiate this claim and show which of these questions matter most, we use the case of a well-understood green transformation policy, the decarbonization of the German housing sector, to illustrate how different closure choices affect final estimates. By doing so, we hope to provide applied researchers with better guidance and tools\footnote{To this end we provide \href{https://github.com/theplatters/OpenCGE}{OpenCGE}, an open-source toolbox that implements all closures used in this paper. Its labor-market, financing and pricing closures are interchangeable modules of one production-network CGE model, so modelers can reproduce our comparisons or apply them to their own input--output tables.} for reflecting on how the diversity of available interpretations of input--output models bears on their problem at hand.

In what follows we will specifically focus on two types of closures, namely labor market closures and financing closures, which can be mapped onto a configuration space with two categorical axes. In this space one axis goes from fundamental scarcity of factors (i.e. fully utilized factor markets) to flexible output adjustment (in the spirit of the Keynesian multiplier), which leads to fundamentally different conceptions of the short run. The second axis focuses on financing and goes from the assumption of immediate refinancing (additional expenditures lead to cuts elsewhere) to assuming some form of external financing that relieves the contemporary economy from immediate refinancing needs. These axes are captured in Section~\ref{sec:two-questions}, which argues that the two poles are domains of applicability rather than rival theories: the output-adjustment pole describes an economy with idle capacity, the scarcity pole one working at its capacity limit, and which of them a study should use is a claim about the state of the world rather than only about modeling technique.
% AG (Oct 2026): added so that R2.3's point (the domain-of-applicability reframing) is in the text, and not only in the
% response letter. Delete the clause and R2.3's response must be reworded accordingly.

The remainder of this paper is structured as follows. Section~\ref{sec:io-bridge} connects transformation research to input--output modeling and introduces the green-investment impulse. Section~\ref{sec:two-questions} develops the two primitive questions -- the nature of the short run and the role of money --, and presents a matrix capturing different scenarios associated with different closure assumptions. Section~\ref{sec:kernel} presents the shared model core inhabited by these different scenarios. Section~\ref{sec:application} runs the program through the matrix, and Section~\ref{sec:beyond-demand-only} extends the demand-only setting by endogenous productivity and markup pricing. Section~\ref{sec:conclusion} summarizes key points, potential follow-ups and relevant limitations of the exercise. Proofs and calibration details are collected in the appendices.
%NOTE: revise at the end


%TODO: We could think once more how to strenghten the contribution – open access codebase would be an idea. in combination with that we could also promote our appendix??????  Idee: Open source package mit dashboard interface, wo man eine Input-output Matrix und einen Shock reinspielt und dann eine Anaylse bekommt, wie die Ergebnisse sich zwischen den Closures unterscheiden. Dazu einfach kernel code abspalten; Kann ich mich dran setzen
